\chapter{REVISÃO DE LITERATURA}

\section{Aprendizado de Máquina no contexto esportivo}

Em 2020, \textit{Horvat \&Job} apresentaram uma revisão abrangente sobre o uso de AM para prever resultados esportivos, 
com foco principal em esportes coletivos, especialmente o futebol. Os autores destacam que o AM é particularmente adequado para esse contexto, 
dada a grande disponibilidade de dados estruturados e não estruturados provenientes de eventos de jogo, sistemas de \textit{tracking} e 
bases históricas consolidadas \cite{Horvat2020}.

Conforme apontado pelos autores, a maior parte dos estudos identificados concentra-se na predição de resultados de partidas individuais, 
especialmente na classificação de vitórias, empates e derrotas. Essa predominância pode ser explicada tanto pelo interesse mercadológico associado 
às apostas esportivas \cite{Abidin2021}, quanto pela própria adequação técnica do problema aos algoritmos supervisionados mais difundidos. 
Entretanto, essa ênfase na previsão de partidas isoladas implica menor atenção a problemas relacionados ao desempenho agregado ao longo de uma temporada, 
como a modelagem da pontuação final de uma equipe a partir da composição de seu elenco.